% Part: computability
% Chapter: computability-theory
% Section: applications-fixed-point

\documentclass[../../../include/open-logic-section]{subfiles}

\begin{document}

\olfileid{cmp}{thy}{apf}
\olsection{تطبيق مبرهنة النقطة الثابتة}

تتيح لنا مبرهنة النقطة الثابتة، في جوهرها، تعريف دوال جزئية قابلة
للحساب بدلالة فهارسها. فمثلًا يمكننا إيجاد فهرس $e$ بحيث، لكل $y$،
\[
\cfind{e}(y) = e + y.
\]
وكمثال آخر يمكن استعمال برهان مبرهنة النقطة الثابتة لتصميم برنامج
بلغة Java أو C++ يطبع نفسه.

تذكر أننا إذا جعلنا، لكل $e$، المجموعة $W_e$ مجال $\cfind{e}$، فإن
المتتالية $W_0$، و$W_1$، و$W_2$،~\dots تعدد المجموعات القابلة للتعداد
حسابيًا. وبعض هذه المجموعات قابل للحساب. وقد نسأل هل توجد خوارزمية
تأخذ قيمة $x$ دخلًا، فإذا صادف أن $W_x$ قابلة للحساب أعادت فهرسًا
لدالتها المميزة. والجواب «لا»: لا توجد خوارزمية كهذه.

\begin{thm}
لا توجد دالة جزئية قابلة للحساب $f$ لها الخاصية الآتية: متى كانت
$W_e$ قابلة للحساب كانت $f(e)$ معرّفة وكانت $\cfind{f(e)}$ دالتها
المميزة.
\end{thm}

\begin{proof}
لتكن $f$ أي دالة جزئية قابلة للحساب؛ سننشئ فهرسًا $e$ بحيث تكون
$W_e$ قابلة للحساب، لكن $\cfind{f(e)}$ ليست دالتها المميزة. باستعمال
مبرهنة النقطة الثابتة يمكننا إيجاد فهرس $e$ بحيث
\[
\cfind{e}(y) \simeq
\begin{cases}
0 & \text{إذا كان $y=0$ وكانت $\cfind{f(e)}(0) \fdefined = 0$} \\
\text{غير معرّفة} & \text{خلاف ذلك.}
\end{cases}
\]
أي إن $e$ ينتج بتطبيق مبرهنة النقطة الثابتة على الدالة المعرّفة
بالمعادلة
\[
g(x,y) \simeq
\begin{cases}
0 & \text{إذا كان $y=0$ وكانت $\cfind{f(x)}(0) \fdefined = 0$} \\
\text{غير معرّفة} & \text{خلاف ذلك.}
\end{cases}
\]
ويمكننا أن نرى بصورة غير صورية أن $g$ جزئية قابلة للحساب كما يأتي:
عند الدخلين $x$ و$y$ تتحقق الخوارزمية أولًا مما إذا كان $y$ يساوي
$0$. فإذا كان كذلك حسبت الخوارزمية $f(x)$، ثم استعملت الآلة الشاملة
لحساب $\cfind{f(x)}(0)$. فإذا توقفت هذه العملية الأخيرة وأعادت~$0$
أعادت الخوارزمية~$0$؛ وإلا لم تتوقف.

لكن لاحظ الآن أنه إذا كانت $f(e)$ غير معرّفة، كان $W_e=\emptyset$
وأخفق شرط تعريف $f(e)$ نفسه. وإذا كانت $f(e)$ معرّفة وكانت
$\cfind{f(e)}(0)$ معرّفة ومساوية لـ$0$، كانت $\cfind{e}(y)$ معرّفة
بالضبط عندما يساوي $y$ العدد $0$، ومن ثم $W_e=\{0\}$. وإذا كانت
$f(e)$ معرّفة، لكن $\cfind{f(e)}(0)$ غير معرّفة أو لا تساوي $0$، كان
$W_e=\emptyset$. وفي الحالتين الأخيرتين لا تكون $\cfind{f(e)}$ الدالة
المميزة لـ~$W_e$، لأنها تعطي الجواب الخطأ عند الدخل $0$.
\end{proof}

\end{document}
